\documentclass[12pt]{article}
\usepackage[utf8]{inputenc}
\usepackage[T1]{fontenc}
\usepackage{hyperref}
\title{Tindi Language Resources for K-12 Curriculum Development}
\author{Compiled by Claude Sonnet 5 (claude-sonnet-5) for Andr\'as Kornai}
\date{}
\begin{document}
\maketitle

\section{Introduction}
Tindi (ISO 639-3 tin; Glottocode tind1238) is an Andic Nakh-Daghestanian language of Dagestan, Russia. The sources seen in search describe it as mainly an oral language, with Avar or Russian used for writing. Speaker counts conflict: roughly 2,150--2,200 in some sources and about 4,500 in the current Wikipedia article. The language is poorly investigated, and search found little.

\section{Primary Educational Resources}
No Tindi primers, textbooks or school programmes turned up in search.

\section{Linguistic and Reference Materials}
\begin{itemize}
\item Tindi--Russian dictionary (\emph{Tindinsko-russkii slovar'}, 2003), seen in a Library of Congress authority record. \url{https://id.loc.gov/authorities/subjects/sh2008005693.html}
\item Omniglot page on Tindi, which mentions a dictionary of about 8,000 words. Its publication details were not confirmed. \url{https://www.omniglot.com/writing/tindi.htm}
\item Wikipedia article on Tindi. It cites Magomedbekova, ``Tindinskij Jazyk'', in \emph{Yazyki mira: Kavkazskie yazyki} (Moscow: Academia, 2001), as a grammar sketch. \url{https://en.wikipedia.org/wiki/Tindi_language}
\item Publications list of Magomed Magomedov (Academia.edu). Its relevance to Tindi was not verified. \url{https://independent.academia.edu/alkaitagi}
\end{itemize}

\section{Religious and Community Resources}
Nothing found in search.

\section{Digital and Audio Resources}
Endangered Languages Project entry, with possible samples. \url{https://www.endangeredlanguages.com/lang/tin?hl=en}

The Estonian Red Book page on Tindi was also found. \url{https://arhiiv.eki.ee/books/redbook/tindis.shtml}

\section{Curriculum Development Context}
Little exists. Avar and Russian are the written languages, and no Tindi orthography or school materials were found. A curriculum would first need a practical orthography, then material built on the 2003 dictionary and the community.
\end{document}